%%%%%%%%%%%%%%%%%%%%%%%%%%%%%%%%%%%%%%%%%%%%%%%%%%%%%%%%%%%%%%%%%%%%%%%%%%%%%%%%
%   Copyright 2009, Oracle and/or its affiliates.
%   All rights reserved.
%
%
%   Use is subject to license terms.
%
%   This distribution may include materials developed by third parties.
%
%%%%%%%%%%%%%%%%%%%%%%%%%%%%%%%%%%%%%%%%%%%%%%%%%%%%%%%%%%%%%%%%%%%%%%%%%%%%%%%%

\chapter{Type Inference}
\chaplabel{type-inference}

A call of a function, a method or an operator whose declaration has static
parameters, or of the constructor of an object that has static parameters,
need not write its static arguments, and most calls do not
(\secref{idn-ref}, \secref{dotted-method-calls}).
The static arguments of a call that does not write them are inferred
statically.
This chapter describes that inference for type parameters (\secref{type-param})
and for \KWD{nat} and \KWD{int} parameters (\secref{natparams}); the inference
of the static arguments of other kinds of static parameter, and the inference
over a whole program of the types that declarations elide
(\secref{type-inference-components}), are not yet described.
\revision{revival-inference}{The Working Draft of February 2011 left this
chapter a heading and two draft notes, the second of which is kept at its end,
although \chapref{multiple-dispatch}, \secref{idn-ref} and
\secref{dotted-method-calls} cite it.
The rule this chapter states is the one the revival decided on 26, 27 and
28~September 2026, and on 29~September 2026 for a lone type parameter with no
narrowest candidate; the compiled type checker was built with this text, and
the cases in which it departs from the text are listed in
\secref{revival-inference}.
The Working Draft had no counterpart to the rule for a numeral whose
conversions tie (\secref{numeral-default}).
A type parameter declared without an \KWD{extends} clause is bounded by
\TYP{Any} (\secref{type-param}); the compiled path still takes \TYP{Object} for
such a parameter of a program's own declarations (row~412 of the revival's gap
ledger).
\chapref{multiple-dispatch} and \chapref{overloaded-declarations} were revised
with this chapter's order, a declaration chosen on the parameter types it
declares and then instantiated (\secref{revival-overloading}).
The team's later Types chapter gives the coercion relation and the valid
instantiations of a quantified type, and says nothing of inference.}

\section{The Static Arguments of a Call}
\seclabel{static-arg-inference}

The static arguments of a call that does not write them are found in two
steps: the declaration that the call invokes is chosen first, and it is
instantiated second.

\emph{Choice.}
The declaration is chosen as \secref{resolving-coercion} describes.
A declaration with static parameters is applicable to a call, without or with
coercion, if some instance of it whose static arguments satisfy their bounds is
applicable to the call.
Declarations are compared on the parameter types they declare, a declaration
with static parameters on its parameter types taken over every instance that
its bounds allow, and not on an instance inferred for the call.
So a coercion is applied only to make a call possible: it never changes which
declaration is chosen when one is applicable to the call without it.

\emph{Instantiation.}
The static arguments of the chosen declaration are then inferred, and with them
the instance of the declaration that the call invokes:
\begin{itemize}
\item A type parameter, or a \KWD{nat} or \KWD{int} parameter, that the
declared type of a parameter mentions other than as the whole type, as
\EXP{b\COLON\TYP{List}\llbracket{}T\rrbracket} mentions \VAR{T}, is fixed by
the argument at that position, whose type must be a subtype of the parameter
type of the instance.
\item Where the call has an expected type (below), a type parameter that the
declared return type mentions other than as the whole type is fixed by the
expected type as well.
\item A type parameter that is the whole declared type of one or more
parameters, as \VAR{T} is of \EXP{x\COLON{}T}, and that is not fixed as above,
takes, among the types of the arguments at those positions and the types to
which they can be coerced, the narrowest under $\morespecificeq$
(\secref{resolving-coercion}) that every one of those arguments is
substitutable for (\secref{applicability-with-coercion}) and that the bounds of
the type parameter permit.
For number arguments this is the narrowest type into which every one of them
converts: \EXP{\mathbb{Z}64} for a \EXP{\mathbb{Z}32} and a
\EXP{\mathbb{Z}64}, and also for a \EXP{\mathbb{Z}32} and an
\EXP{\mathbb{N}32}, though neither of them is a \EXP{\mathbb{Z}64}
(\secref{basic-integers}).
Where there is no such narrowest type, and the rule for numerals
(\secref{numeral-default}) does not settle the tie, the type parameter takes
the intersection of its upper bounds: its declared bound, \TYP{Any} if it has
none (\secref{type-param}), and what the expected type (below), where the call
has one, requires of the declared return type, provided the types of those
arguments are subtypes of it; none of those arguments is converted.
\revision{revival-inference}{The revival's text of 28~September 2026 gave such
a type parameter the union of the types of those arguments; on 29~September 2026
the revival took the rule by which the dispatcher of the type group's paper of
2019 instantiates a type parameter, which gives the bound.}
\item A type parameter that neither an argument nor the expected type fixes as
above, such as a type parameter that only the return type mentions, as its whole
type, likewise takes the intersection of its upper bounds: its declared bound,
\TYP{Any} if it has none (\secref{type-param}), and what the expected type
(below), where the call has one, requires of the declared return type.
Inference never instantiates a static parameter at \TYP{BottomType}, which
programmers must not write (\secref{bottom-type}).
So, for the declaration
\EXP{\VAR{stop}\llbracket{}T\rrbracket(s\COLON\TYP{String})\COLON{}T}, the
variable declaration
\EXP{x\COLON\mathbb{Z}32 = \VAR{stop}(\hbox{\rm``\STR{no}''})} instantiates
\VAR{stop} at \EXP{\mathbb{Z}32}, and the same call with no expected type
instantiates it at \TYP{Any}.
\revision{revival-inference}{The revival's text of 28~September 2026 did not
describe such a type parameter and left open whether inference may give it
\TYP{BottomType}, which the team's draft note at the end of this chapter asks.
The revival decided on 29~September 2026 that the rule by which the dispatcher of
the type group's paper of 2019 instantiates a type parameter that the arguments
do not fix, the intersection of its upper bounds, which never instantiates a
type parameter at \TYP{BottomType}, is the rule of static inference as well.}
\item Every argument is then admitted against the corresponding parameter type
of the instance by substitutability: an argument whose type is a subtype of it
is passed as it is, and any other argument is coerced to it
(\secref{coercion}), as an argument is coerced to the declared type of a
parameter of a declaration without static parameters.
\end{itemize}
The coercions that the instance needs are inserted into the call statically,
as is any coercion that \secref{resolving-coercion} chooses.

\emph{The expected type.}
A call has an expected type in these contexts: as the right-hand side of a
variable or field declaration that declares its type, that type; as the
right-hand side of an assignment to a variable whose type is declared, that
type; as the body of a functional that declares its return type, that type;
as the last expression of a block that has an expected type, where no local
declaration precedes it in the block, the block's expected type; and as a
branch of an \KWD{if} expression with an \KWD{else} clause that has an
expected type, the expected type of the \KWD{if} expression.
The expected type is used whether the call is written by tight juxtaposition,
as \EXP{f(x)} (\secref{juxtameaning}), or as a method invocation or an operator
application.
The return type of the instance that the steps above find must be a subtype of
the expected type.
If they find no such instance, the static arguments are inferred again without
the expected type, and the call is accepted only if the return type of the instance
then found is a subtype of the expected type or can be coerced to it; in the
second case the result of the call is coerced to the expected type
(\secref{coercion}).
So, for the declarations
\EXP{\VAR{wrap}\llbracket{}T\rrbracket(x\COLON{}T)\COLON\TYP{List}\llbracket{}T\rrbracket}
and \EXP{\VAR{id}\llbracket{}T\rrbracket(x\COLON{}T)\COLON{}T},
the variable declaration \EXP{a\COLON\TYP{List}\llbracket\mathbb{Z}64\rrbracket = \VAR{wrap}(3)}
instantiates \VAR{wrap} at \EXP{\mathbb{Z}64}, which the expected type fixes,
and coerces the numeral, while \EXP{b\COLON\mathbb{Z}64 = \VAR{id}(3)}
instantiates \VAR{id} at the numeral's own type (\secref{literals}) and coerces
the result.

\emph{Dispatch.}
The call, with the coercions of its instance inserted, dispatches at run time
as \secref{resolving-overloading} describes, by subtyping and without further
coercion, to the most specific declaration applicable to the converted values.

For example, of the declarations
\EXP{\VAR{op}\llbracket{}T \KWD{extends} \TYP{Number}\rrbracket(a\COLON{}T, b\COLON{}T)}
and \EXP{\VAR{op}(a\COLON\TYP{Any}, b\COLON\TYP{Any})}, the call
\EXP{\VAR{op}(z, w)}, where \VAR{z} has type \EXP{\mathbb{Z}32} and \VAR{w} has
type \EXP{\mathbb{Z}64}, chooses the first, which is applicable without
coercion, since its instance at \TYP{Number} is, and is more specific than the
second; the first is instantiated at \EXP{\mathbb{Z}64}, \VAR{z} is coerced,
and the converted call runs that instance.
\revision{revival-inference}{The interpreter has no static types and applies
the same rule at dispatch, from the run-time types of the arguments: it infers
the static arguments of a declaration with static parameters with coercion,
and instantiates the declaration that its dispatch chooses by the rule for a
type parameter that is a parameter's whole type.
So where the declared type of a variable is wider than the type of its value,
or an argument's static type is \TYP{Any}, it may run a declaration at a
narrower instance than the compiled path, which instantiates it at the static
types.
Where neither an argument's type nor a type into which the arguments convert is
the narrowest, it gives a type parameter that is the whole type of parameters
its declared bound, \TYP{Any} if it has none, as this chapter does: for a
\EXP{\mathbb{Z}64} and an \EXP{\mathbb{R}64}, \TYP{Number} where that is the
bound and \TYP{Any} where the type parameter is declared without an
\KWD{extends} clause.
It does the same where the arguments that fix a type parameter inside other
types, as the results of two function arguments do, have several minimal common
supertypes.
Neither holds of a type parameter whose bound mentions a static parameter,
itself or another, or that is declared with several bounds that the interpreter
cannot meet, as it cannot meet two bounds neither of which is a subtype of the
other: such a type parameter it binds, as before, to the narrowest named common
supertype of the arguments' run-time types that the bounds permit, where there
is one, and where there is none a call that fixes a type parameter of the
second kind fails (row~591 of the revival's gap ledger).
When its dispatch chooses among declarations, it passes over a declaration with
static parameters whose arguments have several minimal common supertypes in
favour of any other declaration applicable to the call without coercion
(row~589 of the revival's gap ledger).
It does not yet fix a type parameter by the expected type.
A type parameter that an argument bounds only from above, as a function
argument bounds the type parameter of its parameter type, it instantiates at the
intersection of the upper bounds that its declared bound and the arguments
place on it: at \EXP{\mathbb{Z}32} for a function of a \EXP{\mathbb{Z}32} where
the type parameter is declared without an \KWD{extends} clause, and at
\TYP{Number} for a function of an \TYP{Any} where it is bounded by
\TYP{Number}; such a type parameter whose bound mentions a
static parameter, or that is declared with several bounds that it cannot meet,
it still instantiates at \TYP{BottomType}; and where the arguments place on
such a type parameter two upper bounds that it cannot meet, as a function over
one trait and a function over another do when neither trait extends the other,
the call fails with an error of unification (row~616 of the revival's gap ledger).
It compares declarations with static parameters on the parameter types they
declare, except generic dotted methods (row~21 of the revival's gap ledger), and
dispatches a converted call again, except a call of an overloaded method.}

\section{A Numeral Whose Conversions Tie}
\seclabel{numeral-default}

A numeral has a type of its own, from which \library\ define coercions to the
number types (\secref{literals}).
Where the types to which a numeral can be converted are incomparable under
$\morespecificeq$, as \EXP{\mathbb{Z}32} and \EXP{\mathbb{N}32} are, each
taking a numeral and neither converting into the other, the numeral is read as
a \EXP{\mathbb{Z}32}, or as a \EXP{\mathbb{Z}64} or a \EXP{\mathbb{Z}} when its
value does not fit in the narrower type.
This holds in either step: in the choice among declarations that are applicable
to the call only with coercion and that the numeral would otherwise leave tied,
and in the instantiation of a type parameter that only numerals fix.
So of the declarations \EXP{g(x\COLON\mathbb{N}32)} and
\EXP{g(x\COLON\mathbb{Z}32)}, the call \EXP{g(0)} invokes the second; and a
type parameter that only numerals fix, whose bounds \EXP{\mathbb{Z}32} and
\EXP{\mathbb{N}32} satisfy and the numeral's own type does not, is
instantiated at \EXP{\mathbb{Z}32} when \EXP{\mathbb{Z}32} holds the numerals'
values.

The rule settles only a tie that a numeral makes.
Where an argument that is not a numeral leaves the declarations applicable to a
call only with coercion tied, no declaration is the most specific, and the call
is a static error.

This chapter does not yet describe:
\begin{itemize}
\item the static arguments of a \KWD{bool}, \KWD{dim}, \KWD{unit} or operator
parameter (\secref{boolparams}, \secref{dimunitparams}, \secref{opr-ident});
\item a \KWD{nat} or \KWD{int} parameter that neither an argument nor the
expected type fixes;
\item the element type of a reduction that only its generator clauses determine
(\secref{reduction-expr}): the rule above gives a static parameter of the big
operator that nothing at the call fixes its bounds, not the type of the
elements;
\item whether an argument whose declared parameter type mentions a static
parameter inside another type may be converted where only a conversion of its
elements would satisfy the bounds of the parameter; under the rule above such
an argument fixes the parameter by subtyping and is not converted;
\item an expected type in any context other than those listed above, such as
an argument of another call, an expression that a local declaration precedes in
a block, an \KWD{if} expression without an \KWD{else} clause, a \KWD{typecase}
expression, the body of a \KWD{for} loop, or a call written by loose
juxtaposition, as \EXP{f\;x};
\item the static arguments of an object with static parameters that is named
without a call, which a program writes (\secref{where-clauses});
\item the inference over a whole program of the types that declarations elide
(\secref{type-inference-components}).
\end{itemize}
\revision{revival-inference}{The compiled type checker instantiates a type
parameter that neither an argument nor the expected type fixes at the
intersection of its upper bounds, as \secref{static-arg-inference} says; on the
compiled path a type parameter of a program's own declaration that is declared
without an \KWD{extends} clause is bounded by \TYP{Object} (row~412 of the
revival's gap ledger), so that a call that fixes such a parameter by nothing
else, where the expected type is one that \TYP{Object} excludes, such as
\TYP{()}, is a static error there.
A reduction's big operator is an argument of the call that the reduction's
desugaring makes, so no expected type reaches it: the compiled type checker
instantiates a generic big operator whose static argument the reduction does
not write at its bound (row~425 of the revival's gap ledger).
The interpreter instantiates a type parameter of a function that nothing at a
call fixes at its bound, \TYP{Any} if it has none, and one whose bound mentions
other static parameters at that bound with their instances, with two
exceptions: a type parameter that only the bound of a type parameter that the
arguments fix mentions, as \EXP{A \KWD{extends} T} mentions \VAR{T}, it
instantiates at the arguments' type, not at its bound, since it has no expected
type to fix it (row~592 of the revival's gap ledger); and a type parameter
declared with several bounds that it cannot meet, as two bounds neither of which
is a subtype of the other, it erases to \TYP{BottomType} (row~591).
A type parameter whose bound mentions the type parameter itself, as the big
operators $\sum$ and $\prod$ declare theirs, and that nothing at a call fixes,
it leaves open: it does not instantiate it, and it admits every value where it
checks a value against that type, as an argument, a variable declared with that
type, a \KWD{typecase} clause or a type assertion.
So a reduction whose element type nothing at the call fixes runs on the types
of its elements, the interpreter's counterpart of the desugaring by the type
\VAR{N} of the expression (\secref{reduction-expr}); but an empty such
reduction takes the identity of \EXP{\mathbb{Z}32} whatever the type of the
expression (row~NEW-W-1 of the revival's gap ledger), and
\EXP{\OPR{BIG}\:\OPR{MINMAX}}, whose reduction joins pairs of elements, stops
(row~NEW-W-2).
It still erases to \TYP{BottomType} the other static parameters of a big
operator that a reduction or a comprehension invokes without its static
arguments, a comprehension's collection being then built at the type of its
elements (row~424 of the revival's gap ledger); it gives a function expression
that declares no return type the return type \TYP{BottomType}, so that a type
parameter that only its result fixes is instantiated there (row~587); and it
leaves a generic method's static parameters uninstantiated (row~21).
The third item of the team's draft note below asks whether inference may give a
static parameter \TYP{BottomType}; the rule of \secref{static-arg-inference},
taken from the type group's paper of 2019, answers that it may not.
Neither implementation infers a \KWD{bool} parameter yet: the compiled type
checker refuses by name to infer a \KWD{bool}, \KWD{dim} or \KWD{unit}
parameter, and the interpreter does not unify a \KWD{bool} one (rows~21
and~307 of the revival's gap ledger).
The example of \secref{boolparams}, a coercion whose \KWD{bool} parameter only
its parameter type mentions, can be applied only if its argument fixes that
parameter.}

\note{ \begin{itemize}
 \item
 There seems to be a circular dependency between inference and juxtaposition
 disambiguation -- you have to know if the subexpressions have arrow types
 to disambiguate, but you can't do that without using inference, which
 requires knowing how the variable of unknown type is used.
\item Type Inference (Dan's email titled ``Comments on a first reading of the spec'' on 06/05/07)
\item Do we want to forbid such cases where type inference infers
\TYP{BottomType}s for static parameters?
 \end{itemize}
}
